% !PRINTED 144
\weberpage{144}
%<<CONT

where the coefficients $A_0,A_1,\ldots,A_m$ depend only on $a$ and on the coefficients still occurring in $S$. We note expressly, however, that if in $S(\alpha,\beta)$ no fractional numerical coefficients occur, then no fractions occur in the coefficients $A_0,A_1,\ldots,A_m$ either.

We now put
\begin{equation}
 \Phi(x)=A_0x^m+A_1x^{m-1}+\cdots+A_{m-1}x+A_m
\tag{4}
\end{equation}
and divide $\Phi(x)$ by $f(x)$, by \S\,3. We obtain a quotient $Q$ and a remainder which, with respect to $x$, is at most of the first degree; thus
\begin{equation}
 \Phi(x)=Qf(x)+A+Bx;
\tag{5}
\end{equation}
Here $A$ and $B$ are now whole functions of $a$ and $b$, and they contain no fractional numerical coefficients if none occur in $S$. If now we put $x=\alpha$, it follows from (5) and (3), since $f(\alpha)$ vanishes, that
\begin{equation}
 S(\alpha,\beta)=A+B\alpha.
\tag{6}
\end{equation}
But since $S(\alpha,\beta)$ and likewise $A,B$ are symmetric, it follows by interchanging $\alpha$ and $\beta$ that
\begin{equation}
 S(\alpha,\beta)=A+B\beta.
\tag{7}
\end{equation}
From this one concludes, since $\alpha$ and $\beta$ are independent variables, that $B=0$, and consequently
\begin{equation}
 S(\alpha,\beta)=A,
\tag{8}
\end{equation}
which proves the fundamental theorem in this case.

\begin{center}
{\large \S.\ 44.}

\vspace{0.5em}
{\large\gesperrt{General Proof of the Principal Theorem.}}
\end{center}

We now suppose that the fundamental theorem has been proved for symmetric functions of \(n-1\) variables, and we derive it, by a procedure entirely analogous to that applied in \S\,43, for \(n\) variables.

Let again
\begin{equation}
 S=S(\alpha_1,\alpha_2,\ldots,\alpha_n)
\tag{1}
\end{equation}
be a whole symmetric function of the \(n\) variables \(\alpha_1,\alpha_2,\ldots,\alpha_n\). If we arrange it according to powers of \(\alpha_1\), and accordingly put
\begin{equation}
 S=S_0\alpha_1^\mu+S_1\alpha_1^{\mu-1}+\cdots+S_{\mu-1}\alpha_1+S_\mu,
\tag{2}
\end{equation}
%CONT>>
